\dykband{EXERCISES IN SWAFL}
{\small Open SWAFL (desktop or online) with the benchmark project. The chapter of the \href{https://github.com/halilunaziru73-creator/Coupled-Water-Fluxes-Pore-Scale-Dynamics-Precision-Irrigation-Mulching-Feedback-Loops/raw/main/Software/v1.1.0/SWAFL_Desktop/docs/manual_v1.1.0/SWAFL_Manual_v1.1.0.pdf}{SWAFL manual} that explains each task is given in brackets; record your values in a table and compare them with what is expected (\emph{Check}).}\par\vspace{3pt}
\begin{exlist}
\item {\color{khblue}[Manual, Chapter 13A]} Open \emph{Drainage design}; record the spacing for the loam and for the clay loam at the same design flow.\par\noindent{\small\itshape Check: the clay loam needs closer drains.}
\item {\color{khblue}[Manual, Chapter 13A]} Change the design flow from 5 to 15 mm\,d$^{-1}$ and plot the spacing.\par\noindent{\small\itshape Check: the spacing falls roughly as $1/\sqrt q$.}
\item {\color{khblue}[Manual, Chapter 13A]} Read the water table through the season and find the wettest period.\par\noindent{\small\itshape Check: the table is highest after the winter rain.}
\item {\color{khblue}[Manual, Chapter 13A]} Read \emph{The same drains in every soil} and rank the soils by the depth of the water table.\par\noindent{\small\itshape Check: the fine soils keep the highest table.}
\end{exlist}
